\subsection{Particularidades del timple}
\label{c:fundamentos_timple}

En su práctica, el timple se caracteriza por dos elementos:

\textbf{En primer lugar}, encontramos dos tipos de rasgueo:
\begin{itemize}
    \item El \textbf{tresillo}~\cite{pedroizquierdo_tresillo}, consistente en tres rasgueos seguidos, normalmente utilizando el índice y el pulgar. La secuencia suele ser: índice abajo, pulgar arriba y índice arriba.
    \item El \textbf{cinquillo}~\cite{pedroizquierdo_cinquillo}, consistente en cinco rasgueos seguidos, normalmente añadiendo a un rasgueo básico el tresillo. La secuencia suele ser: índice abajo, índice arriba, índice abajo, pulgar arriba, índice arriba.
\end{itemize}

\textbf{En segundo lugar}, dentro del folclore canario las \textbf{tonalidades}~\cite{timplismo_tonalidades}\cite{raquel_tonalidades_video} reciben un nombre propio \tabref{tab:tonalidades_timple}, siendo los grados I, IV, V7 los mas utilizados.

\begin{table}[H]
    \centering
    \scriptsize
    \begin{tabular}{l c c c}
    \hline
    \textbf{Nombre} & \textbf{I} & \textbf{IV} & \textbf{V7} \\
    \hline
    Tendío       & $DoM$ & $FaM$ & $Sol7$ \\
    Por el cinco & $ReM$ & $SolM$ & $La7$ \\
    Por el nueve & $MiM$ & LaM & Si7 \\
    La zorra     & $FaM$ & $SibM$ & $Do7$ \\
    Por el uno   & $SolM$ & $DoM$ & $Re7$ \\
    Por el seis  & $LaM$ & $ReM$ & $Mi7$ \\
    \hline
    \end{tabular}
    \caption{Tonalidades, grados y acordes con sus nombres folcloricos}\label{tab:tonalidades_timple}
\end{table}

\textbf{Además}, existe una variante del timple conocido como \textbf{contra} con afinación \textit{La, Mi, Si, Sol y Re}